%% ============================================================
%%  TECT Research Programme: Project Roadmap
%%  Math01 -- Math36 Structured into Six Research Projects
%%  Author: Jusang Lee
%%  Date: April 2026
%% ============================================================

\documentclass[reprint,aps,pre,onecolumn,superscriptaddress,amsmath,amssymb]{revtex4-2}

\usepackage{graphicx}
\usepackage{dcolumn}
\usepackage{bm}
\usepackage{hyperref}
\usepackage{braket}
\usepackage{amsthm}
\usepackage{array}
\usepackage{longtable}
\usepackage{booktabs}
\usepackage{xcolor}
\usepackage{mdframed}

%% ---- TECT standard macros ----
\newcommand{\Ltect}{\mathcal{L}_{\mathrm{TECT}}}
\newcommand{\mstar}{m^{*}}
\newcommand{\Dmu}{D_{\mu}}
\newcommand{\Nloop}{N_{\mathrm{loop}}}
\newcommand{\HilD}{\mathcal{H}_{D}}
\newcommand{\HilL}{\mathcal{H}_{L}}
\newcommand{\Gstar}{G_{*}}
\newcommand{\Psifield}{\Psi(x)\in\mathbb{C}^{3}}

\newtheorem{theorem}{Theorem}[section]
\newtheorem{lemma}[theorem]{Lemma}
\newtheorem{corollary}[theorem]{Corollary}
\newtheorem{proposition}[theorem]{Proposition}
\newtheorem{openproblem}[theorem]{Open Problem}

\begin{document}

\title{Topological Energy Condensate Theory (TECT):\\
  Research Programme Roadmap --- Projects I through VI\\
  (Documents Math01--Math36)}

\author{Jusang Lee}
\email{jusanglee@kakao.com}
\affiliation{Independent Researcher, Incheon, Korea}

\date{April 2026}

\begin{abstract}
Topological Energy Condensate Theory (TECT) derives gravity, gauge fields, and Standard
Model fermions from a single underlying $\mathbb{C}^{3}$-valued energy condensate whose
equilibrium configuration is Body-Centered Cubic (BCC).
The research programme documented in Math01--Math36 decomposes naturally into six
sequential projects:
(I) BCC ground-state theorem via $\mathbb{Z}_{2}$-symmetric Landau--Brazovskii field
theory ($\Nloop = 6$ for BCC versus $2$ for FCC and $0$ for SC);
(II) emergent $U(1)$ gauge structure from transverse BCC fluctuations;
(III) topological fermion emergence via Bloch-shell Weyl nodes and mass-protection;
(IV) microscopic Dirac carrier acceptance (Clifford closure, Gates 2--4);
(V) three-generation flavor structure and the SM$+$GR infrared limit;
(VI) full numerical PDE certification on the Class~II working branch with parameters
$(\mu^{2},\lambda,\gamma)=(0.26,-0.43,1.62)$.
Each project is self-contained, carries its own theorem-level deliverables, and feeds
sequentially into the next.
\end{abstract}

\maketitle

%% ============================================================
\section{Overview: The TECT Logical Chain}
\label{sec:overview}
%% ============================================================

The fundamental hypothesis of TECT is:

\begin{equation}
  \boxed{
    \text{Universe} = \text{BCC energy condensate } \Psi_{0}
    \;\xrightarrow{\;\text{fluctuations}\;}\;
    \text{gauge fields} + \text{fermions} + \text{gravity}.
  }
  \label{eq:master-chain}
\end{equation}

The condensate field is a complex vector order parameter
$\Psifield$, evolving under the microscopic action:

\begin{equation}
  \mathcal{L}_{\mathrm{mic}}
  =
  \mathcal{L}_{\Psi}^{\mathrm{core}}[\Psi]
  + \frac{M_{X}^{2}}{2}X_{i}^{a}X_{i}^{a}
  + \alpha_{X} X_{i}^{a} J_{i}^{a}[\Psi]
  + \beta_{X}  X_{i}^{a} K_{i}^{a}[\Psi],
  \label{eq:mic-action}
\end{equation}

where $X_{i}^{a}$ is a heavy mediator field (Class~II branch) and the core
nonlinear operator is:

\begin{equation}
  \mathcal{L}_{\Psi}^{\mathrm{core}}
  =
  -\mu^{2}\Psi + Z\nabla^{2}\Psi - Y\nabla^{4}\Psi
  - \lambda|\Psi|^{2}\Psi - \gamma|\Psi|^{4}\Psi.
  \label{eq:core-op}
\end{equation}

The six research projects map onto the logical steps of
Eq.~\eqref{eq:master-chain} as follows:

\begin{center}
\renewcommand{\arraystretch}{1.4}
\begin{tabular}{clll}
\toprule
\textbf{Project} & \textbf{Files} & \textbf{Physical Objective} & \textbf{Status} \\
\midrule
I   & Math01--05 & BCC ground-state uniqueness      & \textcolor{green!60!black}{Proved} \\
II  & Math06--09 & Emergent $U(1)$ gauge structure   & \textcolor{green!60!black}{Proved} \\
III & Math10--14 & Topological Weyl/Dirac fermions   & \textcolor{green!60!black}{Proved} \\
IV  & Math15--24 & Microscopic Dirac carrier audit   & \textcolor{orange}{Certificate pending} \\
V   & Math25--30 & Flavor structure \& SM+GR limit   & \textcolor{orange}{Framework complete} \\
VI  & Math31--36 & Numerical PDE certification       & \textcolor{red!70!black}{Active frontier} \\
\bottomrule
\end{tabular}
\end{center}

%% ============================================================
\section{Project I: BCC Ground-State Theorem}
\label{sec:proj1}
%% ============================================================
\textit{Documents: Math01, Math02, Math03, Math04, Math05}

\subsection{Core Objective}

Prove that the BCC Bravais lattice is the \emph{unique} global minimum of the
$\mathbb{Z}_{2}$-symmetric Landau--Brazovskii free energy under pure quartic
(4-wave) fluctuations at the 1-loop level.

\subsection{Governing Free Energy}

The effective free energy functional restricted to a single $k_{0}$-shell is:
\begin{align}
  F[\{\Psi_{\mathbf{k}}\}]
  &=
  \sum_{|\mathbf{k}|=k_{0}} \left[ r + \kappa(k^{2}-k_{0}^{2})^{2} \right] |\Psi_{\mathbf{k}}|^{2}
  \notag\\
  &\quad
  + u \sum_{\substack{\mathbf{k}_{1}+\mathbf{k}_{2}+\mathbf{k}_{3}+\mathbf{k}_{4}=0\\
    |\mathbf{k}_{i}|=k_{0}}}
  \Psi_{\mathbf{k}_{1}}\Psi_{\mathbf{k}_{2}}\Psi_{\mathbf{k}_{3}}\Psi_{\mathbf{k}_{4}},
  \label{eq:LB-free-energy}
\end{align}
with $\mathbb{Z}_{2}$ symmetry ($\Psi\to -\Psi$) forbidding all cubic invariants.

\subsection{Main Result}

\begin{theorem}[BCC Selection Theorem --- Math01]
  \label{thm:BCC}
  Under pure quartic $\mathbb{Z}_{2}$-symmetric fluctuations on a single isotropic
  shell at $k_{0}$, the non-trivial 4-wave resonance loop multiplicities are:
  \begin{equation}
    \Nloop^{\mathrm{BCC}} = 6,
    \qquad
    \Nloop^{\mathrm{FCC}} = 2,
    \qquad
    \Nloop^{\mathrm{SC}} = 0.
    \label{eq:Nloop}
  \end{equation}
  At 1-loop, the fluctuation-induced free-energy shift satisfies
  $\Delta F_{\mathrm{BCC}} < \Delta F_{\mathrm{FCC}} < \Delta F_{\mathrm{SC}}$,
  and the BCC phase is selected as the global minimum via a
  fluctuation-induced first-order transition.
\end{theorem}

\subsection{Robustness Proofs (Math02--05)}

\begin{description}
  \item[Math02 --- Universality A:]
    Finite shell thickness $\Delta k \neq 0$ leaves the inequality
    $\Nloop^{\mathrm{BCC}} > \Nloop^{\mathrm{FCC}}$ intact.
    The isotropic broad-shell kernel $G^{-1}(q) = r + \kappa(q^{2}-k_{0}^{2})^{2}$
    is treated perturbatively; the BCC selection criterion survives at all
    $\Delta k / k_{0} \ll 1$.

  \item[Math03 --- Universality B:]
    The equal-amplitude ansatz $A_{i} = A/\sqrt{N_{S}}$ is not an assumption
    but a consequence. Allowing amplitude variation
    $A_{i} = \bar{A} + \delta A_{i}$ with $\sum\delta A_{i} = 0$ shows the
    equal-amplitude BCC configuration is a strict local minimum of the
    quartic free energy.

  \item[Math04--05 --- Multi-Shell Extension:]
    The resonance network is generalized to wavevectors distributed across
    $K$ shells, $S = \bigcup_{a=1}^{K} S_{a}$, each satisfying antipodal
    symmetry. The BCC selection theorem extends to the multi-shell case
    when the dominant shell satisfies the single-shell dominance condition.
\end{description}

\subsection{Deliverables}
\begin{enumerate}
  \item Exact combinatorial count $\Nloop$ for BCC, FCC, SC (complete).
  \item 1-loop free-energy comparison $\Delta F_{\mathrm{BCC}} < \Delta F_{\mathrm{FCC}}$ (complete).
  \item Universality under finite $\Delta k$ and amplitude variation (complete).
  \item Formal connection to Wigner-Seitz 14-faced truncated octahedron (complete).
\end{enumerate}

%% ============================================================
\section{Project II: Emergent Gauge Structure}
\label{sec:proj2}
%% ============================================================
\textit{Documents: Math06, Math07, Math08, Math09}

\subsection{Core Objective}

Derive emergent $U(1)$ and higher gauge structures from low-energy fluctuations
around the BCC-condensed background $\Psi_{0}$.

\subsection{Linearization around the BCC Background}

Expanding $\Psi = \Psi_{0} + \delta\Psi$ around the stationary BCC point:
\begin{equation}
  S[\Psi_{0}+\delta\Psi]
  = S[\Psi_{0}]
  + \frac{1}{2}\,\delta^{2}S[\Psi_{0};\delta\Psi,\delta\Psi]
  + O(\delta\Psi^{3}).
  \label{eq:background-expand}
\end{equation}
Since $\Psi_{0}$ is a stationary point, the linear term vanishes.
For each wavevector $\mathbf{k}$, the fluctuation decomposes as:
\begin{equation}
  \delta\mathbf{u}
  = u_{L}\,\hat{\mathbf{k}}
  + u_{1}\,\mathbf{e}_{1}
  + u_{2}\,\mathbf{e}_{2}.
  \label{eq:fluctuation-decomp}
\end{equation}

\subsection{Main Results}

\begin{theorem}[$U(1)$ Emergence --- Math06, Math07]
  \label{thm:U1}
  In the infrared limit $|\mathbf{k}|a \ll 1$, the discrete anisotropy of the
  BCC lattice is suppressed by higher-order derivative corrections. The
  transverse sector $\{u_{1}, u_{2}\}$ carries an internal
  $SO(2) \simeq U(1)$ symmetry on the transverse polarization plane.
  The emergent gauge-covariant derivative takes the form:
  \begin{equation}
    \Dmu \sim \partial_{\mu} + i A_{\mu},
    \qquad A_{\mu} \in \mathfrak{u}(1),
    \label{eq:U1-covariant}
  \end{equation}
  where $A_{\mu}$ is identified with the photon field.
\end{theorem}

\begin{theorem}[Global $U(1)$ Bundle --- Math07]
  \label{thm:CP2-bundle}
  The low-energy orientational degree of freedom $z(x)\in\mathbb{C}^{3}$
  with $z^{\dagger}z = 1$ defines a principal $U(1)$ bundle over $\mathbb{CP}^{2}$.
  The corresponding Berry connection gives a geometric $U(1)$ gauge field
  with curvature $F = dA$.
\end{theorem}

\begin{theorem}[Shell-Protected Dirac Couplings Non-Vanishing --- Math09]
  \label{thm:Dirac-coupling}
  For any non-zero shell wavevector $\Gstar \in \mathbb{R}^{3}\setminus\{0\}$,
  the reduced shell coefficients of the symmetry-allowed microscopic enriched
  tensors $\mathsf{N}_{ij}^{a}$ are generically non-vanishing:
  \begin{equation}
    v_{\alpha} = e_{i}^{\alpha}\,\mathsf{N}_{ij}^{a}\,\Gstar_{j} \neq 0,
    \qquad \alpha = 1, 2.
    \label{eq:shell-coupling}
  \end{equation}
  This guarantees a non-trivial linear coupling between the
  transverse modes and the condensate background.
\end{theorem}

\subsection{Deliverables}
\begin{enumerate}
  \item Emergent $U(1)$ gauge field from transverse BCC fluctuations (complete).
  \item $\mathbb{CP}^{2}$ geometric bundle structure (complete).
  \item Non-vanishing Dirac coupling theorem (complete).
  \item Candidate $SU(2)\times U(1)$ structure from the full $\mathbb{C}^{3}$ bundle (in progress).
\end{enumerate}

%% ============================================================
\section{Project III: Topological Fermion Emergence}
\label{sec:proj3}
%% ============================================================
\textit{Documents: Math10, Math11, Math12, Math13, Math14}

\subsection{Core Objective}

Establish that massless Weyl/Dirac fermions emerge at the Bloch-band crossings
on the BCC reciprocal shell, protected by residual valley symmetry and little-group
representation theory.

\subsection{Setup: Bloch Operator and Valley Structure}

The microscopic BCC linearization yields a Bloch Hamiltonian $H(\mathbf{p})$
whose spectrum has crossings at selected shell wavevectors $\pm\Gstar$.
The low-energy doubled block near such a crossing is:
\begin{equation}
  H_{D}(\mathbf{p})
  =
  \begin{pmatrix}
    h_{+}(\mathbf{p}) & 0 \\
    0 & h_{-}(\mathbf{p})
  \end{pmatrix},
  \qquad
  h_{-}(\mathbf{p}) = -h_{+}(-\mathbf{p}),
  \label{eq:valley-pair}
\end{equation}
where $h_{\pm}$ acts on the two-component Weyl doublet at each valley $\pm\Gstar$.

\subsection{Main Results}

\begin{theorem}[Weyl Node Emergence --- Math10, Math11]
  \label{thm:Weyl}
  After a shell-adapted linear change of momentum coordinates, the linearized
  Bloch crossing takes the standard Weyl-paired form:
  \begin{equation}
    H_{D}(\mathbf{p})
    = \sum_{i=1}^{3} v_{i}\,p_{i}\,\alpha_{i},
    \qquad
    \alpha_{i} := \tau_{z} \otimes \sigma_{i},
    \label{eq:Weyl-form}
  \end{equation}
  where $\sigma_{i}$ acts on the Weyl doublet and $\tau_{i}$ acts on the valley pair
  $(+\Gstar, -\Gstar)$.
\end{theorem}

\begin{theorem}[Mass Protection --- Math12]
  \label{thm:mass-protect}
  The residual valley symmetry $\mathbb{Z}_{2}^{\mathrm{valley}}$ acting as
  $\tau_{z}$ forbids all constant Hermitian operators of the form $m\,\tau_{0}\otimes\sigma_{0}$
  from appearing in $H_{D}(\mathbf{p})$.
  Consequently, the Weyl crossing at $p = 0$ is topologically protected:
  \begin{equation}
    \det H_{D}(0) = 0,
    \label{eq:mass-zero}
  \end{equation}
  and no perturbation consistent with the residual symmetry can open a mass gap.
\end{theorem}

\begin{theorem}[Little-Group Mismatch Decoupling --- Math13]
  \label{thm:LG-mismatch}
  Let $\mathcal{H}_{\mathrm{full}} = \mathcal{H}_{D} \oplus \mathcal{H}_{R}$
  be the full Hilbert space decomposed into shell-relevant ($D$) and remote ($R$) sectors.
  If $\mathcal{H}_{D}$ and $\mathcal{H}_{R}$ carry inequivalent irreducible
  representations of the little group $H_{*}$, then the constant off-diagonal
  mixing block satisfies:
  \begin{equation}
    V(0) = 0,
    \label{eq:V-zero}
  \end{equation}
  establishing that remote-sector mixing cannot contaminate the Dirac block at $p=0$.
\end{theorem}

\begin{theorem}[Clifford Algebra Dimension --- Math14]
  \label{thm:Clifford}
  The reduced Branch~B symbol
  $H_{B}(\mathbf{p}) = v_{\parallel}p_{\parallel}\sigma_{3} + v_{\perp}(p_{1}\sigma_{1}+p_{2}\sigma_{2})$
  acts on the tangent fermion $\mathbb{C}^{2}_{\mathrm{int}}$.
  Combined with the valley doublet, the full low-energy Hilbert space is
  $\mathcal{H}_{\mathrm{low}} = \mathbb{C}^{2}_{\mathrm{int}} \otimes \mathbb{C}^{2}_{\mathrm{valley}}$,
  realizing a minimal $4\times 4$ Dirac representation.
\end{theorem}

\subsection{Deliverables}
\begin{enumerate}
  \item Weyl node at BCC shell crossing (complete).
  \item Mass protection theorem via $\mathbb{Z}_{2}^{\mathrm{valley}}$ (complete).
  \item Little-group mismatch decoupling (complete).
  \item Clifford closure on $\mathbb{C}^{4}$ Dirac representation (complete).
\end{enumerate}

%% ============================================================
\section{Project IV: Microscopic Dirac Carrier Acceptance}
\label{sec:proj4}
%% ============================================================
\textit{Documents: Math15, Math16, Math17, Math18, Math19, Math20, Math21, Math22, Math23, Math24}

\subsection{Core Objective}

Perform the finite microscopic operator audit confirming that the BCC linearized
Bloch operator admits a Dirac carrier satisfying all four acceptance gates:
\begin{equation}
  \boxed{
    \text{Gate 1: carrier isolation}
    \;\wedge\;
    \text{Gates 2--3: spectrum audit}
    \;\wedge\;
    \text{Gate 4: mass protection.}
  }
  \label{eq:gates}
\end{equation}

\subsection{Programme Structure}

\paragraph{Math15 --- Finite Audit Reformulation.}
The open microscopic questions are recast as a finite operator audit on the
actual linearized Bloch operator around the refined locked TECT branch.
The projector onto the condensate orientation is:
\begin{equation}
  P_{0} = z_{0}z_{0}^{\dagger},
  \qquad \Psi_{0} = \phi_{0} z_{0},
  \qquad z_{0} \in \mathbb{C}^{3},\; |z_{0}|=1.
  \label{eq:projector}
\end{equation}

\paragraph{Math16--17 --- Honest PDE Status.}
The current derivation operates on the strongest explicit working microscopic branch
$\mathcal{L}_{\mathrm{mic}}$ (Eq.~\ref{eq:mic-action}) rather than the full
unconditional action. The PDE status is recorded as:
\begin{equation}
  \boxed{
    \text{proof chain complete};\quad
    \text{microscopic coefficient certificate: pending.}
  }
  \label{eq:PDE-status}
\end{equation}

\paragraph{Math18 --- Clifford Closure.}
The reduced symbol on the doubled valley space yields the
$4\times 4$ Dirac Hamiltonian. The Clifford algebra
$\mathrm{Cl}(3,0)$ is shown to close on the four generators
$\{\tau_{z}\otimes\sigma_{i}\}_{i=1}^{3}$ plus $\tau_{x}\otimes\sigma_{0}$.

\paragraph{Math20 --- Explicit Enriched Shell Basis.}
The 4-dimensional Dirac-carrier candidate is constructed explicitly:
\begin{equation}
  \mathcal{H}_{D}^{(N)}
  \cong
  \mathrm{span}\{|\Gstar,\sigma,\tau\rangle\}_{\sigma\in\{+,-\},\,\tau\in\{+,-\}}.
  \label{eq:HD-explicit}
\end{equation}

\paragraph{Math21 --- Carrier Acceptance Stage.}
The frontier is identified as:
\begin{equation}
  \boxed{
    \text{carrier acceptance (Gates 2--3)} + \text{mass protection (Gate 4).}
  }
  \label{eq:frontier}
\end{equation}

\paragraph{Math22--23 --- Same-Shell Non-Enriched Benchmark.}
The benchmark Hamiltonian $H_{\mathrm{non\text{-}enr}}^{(N)}(0)$ is computed
for the minimal truncation $N=1$. Spectral non-vanishing at $p=0$ is verified:
$\det H_{\mathrm{non\text{-}enr}}^{(1)}(0) \neq 0$.

\paragraph{Math24 --- Small-$p$ Drift Control.}
The perturbation $H^{(N)}(p) - H^{(N)}(0)$ is bounded uniformly on $|p| \leq \rho$
by a Lipschitz estimate:
\begin{equation}
  \bigl\| H^{(N)}(p) - H^{(N)}(0) \bigr\| \leq C_{\mathrm{drift}} \cdot |p|,
  \qquad |p| \leq \rho,
  \label{eq:drift-bound}
\end{equation}
closing the small-$p$ stability argument for the carrier spectrum.

\subsection{Deliverables}
\begin{enumerate}
  \item Gate 1 (carrier isolation): complete.
  \item Gates 2--3 (spectrum non-vanishing at $p=0$, Lipschitz bound): complete.
  \item Gate 4 (mass protection): complete.
  \item Full certificate with explicit microscopic coefficients: \emph{pending --- Project VI.}
\end{enumerate}

%% ============================================================
\section{Project V: Three-Generation Flavor Structure and SM+GR Infrared Limit}
\label{sec:proj5}
%% ============================================================
\textit{Documents: Math25, Math26, Math27, Math28, Math29, Math30}

\subsection{Core Objective}

Establish the three-generation left-handed family space $\mathcal{H}_{L}$
from the BCC condensate structure, and formulate the minimal theorem
checklist for an SM$+$GR-type infrared description.

\subsection{Flavor Branch: Three-Generation Emergence}

\paragraph{Math25 --- Local Uniqueness of the Flavor Branch.}
The phenomenologically successful TECT branch is shown to be a locally
unique stationary point of the reduced free energy on the flavor sector.
The $3\times 3$ Hessian of the reduced potential is computed and verified
to have positive eigenvalues at the selected branch point.

\paragraph{Math26 --- SM+GR Infrared Checklist.}
The main infrared target theorem is:
\begin{equation}
  \boxed{
    \text{TECT}_{\mathrm{IR}} \supseteq \text{SM} + \text{GR},
  }
  \label{eq:IR-limit}
\end{equation}
meaning that the long-distance effective description of TECT
must contain the Standard Model gauge group $SU(3)\times SU(2)\times U(1)$
minimally coupled to the metric $g_{\mu\nu}$.
The minimal theorem checklist enumerates 7 required lemmas.

\paragraph{Math27 --- Left-Handed Family Space Isolation (K1).}
The projector onto the left-handed sector yields:
\begin{equation}
  \mathcal{H}_{L} = \mathrm{Ran}(P_{L}),
  \qquad
  \dim \mathcal{H}_{L} = 3.
  \label{eq:HL-dim}
\end{equation}
This 3-dimensional space is identified with the three quark/lepton generations.

\paragraph{Math28 --- Kinetic Gram Matrix (K2A).}
On the common left-handed family space $\mathcal{H}_{L}$, the kinetic
Gram matrix is:
\begin{equation}
  G_{IJ} = \langle L_{I} | \hat{K} | L_{J} \rangle,
  \qquad I, J = 1, 2, 3,
  \label{eq:Gram}
\end{equation}
where $\hat{K}$ is the sector-reduced low-energy kinetic operator.
The matrix $G_{IJ}$ is Hermitian and must be shown positive-definite for
physical normalizability of all three generations.

\paragraph{Math29 --- Full Field Definition.}
The TECT condensate field is fully specified as
$\Psi(x)\in\mathbb{C}^{3}$, with the BCC ground state realized as a
sum over the 12-vector reciprocal star:
\begin{equation}
  \Psi_{0}(\mathbf{x})
  = \phi_{0} \sum_{\mathbf{G}\in\mathcal{S}_{\mathrm{BCC}}} z_{\mathbf{G}}\,
    e^{i\mathbf{G}\cdot\mathbf{x}},
  \quad
  |\mathcal{S}_{\mathrm{BCC}}| = 12.
  \label{eq:BCC-ground-state}
\end{equation}

\paragraph{Math30 --- PDE Completion Protocol.}
The carrier-level certificate is established; the remaining task is
to populate it with actual branch data from the full TECT linearization:
\begin{equation}
  \boxed{
    \text{proof complete};\quad
    \text{certificate incomplete (finite microscopic audit required).}
  }
  \label{eq:completion-status}
\end{equation}

\subsection{Deliverables}
\begin{enumerate}
  \item $\dim\mathcal{H}_{L} = 3$ (three generations): complete.
  \item Kinetic Gram matrix $G_{IJ}$ positive-definite: framework complete,
        explicit values pending Project VI.
  \item SM$+$GR infrared theorem checklist (7 lemmas): framework complete.
  \item BCC ground-state field decomposition: complete.
\end{enumerate}

%% ============================================================
\section{Project VI: Numerical PDE Certification}
\label{sec:proj6}
%% ============================================================
\textit{Documents: Math31, Math32, Math33, Math34, Math35, Math36}

\subsection{Core Objective}

Complete the finite microscopic audit and numerically certify the PDE
coefficient extraction for the Class~II working branch, anchoring all
physical parameters.

\subsection{Operational Parameters (Locked)}

The operational backend is frozen at:
\begin{equation}
  \boxed{
    (\mu^{2},\,\lambda,\,\gamma) = (0.26,\;-0.43,\;1.62).
  }
  \label{eq:locked-params}
\end{equation}

\subsection{Sequential Steps}

\paragraph{Math31 --- Enlarged Quartic Reduced Shell Kernel.}
Going beyond the failed $(u,v,\kappa)$ truncation, the next quartic reduced
kernel is constructed with the full four-class structure. The corrected seed profile
at locked parameters~\eqref{eq:locked-params} is used as the expansion point.

\paragraph{Math32 --- Residual Quartic Closure.}
The exact four-class upgrade of the quartic closure is performed, correcting
systematic errors in the $(u,v,\kappa)$ truncation:
\begin{equation}
  F_{4} = u_{1}\rho_{1} + u_{2}\rho_{2} + u_{3}\rho_{3} + u_{4}\rho_{4},
  \label{eq:quartic-closure}
\end{equation}
where $\rho_{1},\ldots,\rho_{4}$ are the four independent quartic invariants
on the BCC reciprocal star.

\paragraph{Math33 --- Extended PDE Coefficient Extraction.}
The PDE coefficient extraction protocol is applied to the corrected quartic
closure to obtain the physical reduced PDE parameters:
\begin{equation}
  \lambda_{\parallel} = \frac{\lambda_{c}}{2}I_{3},
  \qquad
  \alpha = \beta = -\frac{\beta_{X}\gamma_{\perp}}{M_{X}^{2}}I_{c}.
  \label{eq:PDE-coeffs}
\end{equation}

\paragraph{Math34 --- Seed Profile Extraction.}
The explicit seed profile is inserted into the reference reduced PDE to extract
the finite Bloch operator matrix elements:
\begin{align}
  \mathcal{U}_{j}(0)
  &=
  \phi_{0}\sum_{\mathbf{G}\in\mathcal{S}} J_{j}(\mathbf{G})\,z_{0},
  \label{eq:seed-extract}
\end{align}
enabling direct computation of the kinetic Gram matrix $G_{IJ}$.

\paragraph{Math35 --- Canonical Locked-Basis Reading.}
The longitudinal PDE coefficient $u_{\parallel}$ is extracted in the canonical
locked basis via the $J$-matrix at the selected intervalley harmonic:
\begin{equation}
  u_{\parallel}
  =
  \frac{1}{\phi_{0}}\,
  \hat{e}_{\parallel}^{\dagger}\,\mathcal{U}(0)\,\hat{e}_{\parallel}.
  \label{eq:u-parallel}
\end{equation}

\paragraph{Math36 --- Class~II Elimination and BCC Linearization.}
The heavy mediator $X_{i}^{a}$ is integrated out exactly via its algebraic
Euler--Lagrange equation:
\begin{equation}
  X_{i}^{a}
  = -\frac{1}{M_{X}^{2}}\bigl(\alpha_{X}J_{i}^{a} + \beta_{X}K_{i}^{a}\bigr),
  \label{eq:X-elim}
\end{equation}
yielding the final effective action from which the full BCC linearization
is performed. This is the active frontier of the TECT programme.

\subsection{Deliverables}
\begin{enumerate}
  \item Enlarged quartic kernel (four-class): in progress.
  \item Explicit PDE coefficient values $(\lambda_{\parallel}, \alpha, \beta)$: in progress.
  \item Positive-definite Gram matrix $G_{IJ}$ with explicit entries: target.
  \item Class~II heavy-field elimination and full BCC linearization: target.
  \item Spectral gap $m^{*} = 0.3138 \pm \delta m^{*}$ derived analytically: target.
\end{enumerate}

%% ============================================================
\section{Summary: Logical Dependency Graph}
\label{sec:summary}
%% ============================================================

The six projects form a strict directed acyclic dependency chain:

\begin{equation}
  \underbrace{\text{I (BCC)}}_{\text{Math01--05}}
  \;\longrightarrow\;
  \underbrace{\text{II (Gauge)}}_{\text{Math06--09}}
  \;\longrightarrow\;
  \underbrace{\text{III (Fermion)}}_{\text{Math10--14}}
  \;\longrightarrow\;
  \underbrace{\text{IV (Carrier)}}_{\text{Math15--24}}
  \;\longrightarrow\;
  \underbrace{\text{V (Flavor)}}_{\text{Math25--30}}
  \;\longrightarrow\;
  \underbrace{\text{VI (PDE)}}_{\text{Math31--36}}.
  \label{eq:dependency-chain}
\end{equation}

The current status of the overall TECT programme is:
\begin{align}
  &\text{Projects I--III}: \text{theorem-level proofs complete.} \notag\\
  &\text{Project IV}: \text{carrier acceptance gates 1--4 proved;
    explicit coefficient certificate pending.} \notag\\
  &\text{Project V}: \text{framework complete;
    explicit Gram matrix entries pending.} \notag\\
  &\text{Project VI}: \text{active frontier.
    Class~II elimination (Math36) is the immediate next step.}
  \label{eq:overall-status}
\end{align}

\subsection{Immediate Next Step: Math37 Target}

Completion of Project~VI (specifically Math36) requires:
\begin{enumerate}
  \item Insert Eq.~\eqref{eq:X-elim} into $\mathcal{L}_{\mathrm{mic}}$
        to obtain the exact induced quartic term.
  \item Perform the BCC Bloch decomposition of the resulting effective Lagrangian.
  \item Extract the reduced PDE matrix coefficients $(\lambda_{\parallel}, v_{\perp}, v_{\parallel})$
        numerically at $(\mu^{2},\lambda,\gamma) = (0.26, -0.43, 1.62)$.
  \item Verify $\det H_{D}(0) = 0$ (Weyl node survives Class~II elimination).
  \item Compute the spectral gap $m^{*}$ and compare to the simulation value
        $m^{*} = 0.3138$.
\end{enumerate}
Successful completion of steps 1--5 constitutes the closing certificate
for Project~VI and the first complete end-to-end analytical derivation
connecting the BCC condensate to a measurable physical parameter.

\end{document}
